% 子图专项测试。42 变体里的子图只覆盖了示例正文里那几张，
% subcapcenterlast 分支、subtable、\subref 一律走不到，换宏包时需要单独看。
%
% 文件里的写法按 subfigpkg 分成两套，用 \ifhitmodern 切换，
% 两套排出来的版面应当一致。
\documentclass[type=master,campus=harbin,fontset=fandol]{hithesisbook}
\hitsetup{info={ctitle={子图测试},cauthor={测试},csupervisor={测试},
  cxueke={工学},csubject={测试},caffil={测试},ckeywords={甲,乙},
  etitle={Subfigure Test},eauthor={T},esupervisor={T},exueke={E},
  esubject={E},eaffil={E},ekeywords={a,b}}}
\begin{document}
\begin{cabstract}摘要。\end{cabstract}
\begin{eabstract}Abstract.\end{eabstract}
\makecover
\listoffigures
\listoftables
\mainmatter
\chapter{子图}

\section{两个子图，单语子图题}
\label{sec:plain}
\begin{figure}[!h]
\setlength{\subfigcapskip}{-1bp}
\centering
\subfigure[左边]{\label{sf:a}\includegraphics[width=0.3\textwidth]{golfer}}
\hspace{2em}
\subfigure[右边]{\label{sf:b}\includegraphics[width=0.3\textwidth]{golfer}}
\bicaption[fig:plain]{}{单语子图题}{Fig.$\!$}{Monolingual subcaptions}
\end{figure}

引用：\ref{sf:a}、\ref{sf:b}、\ref{fig:plain}。

\section{双语子图题}
\begin{figure}[!h]
\setlength{\subfigcapskip}{-1bp}
\centering
\begin{minipage}{\textwidth}
\centering
\subfigure[\bisubcaption{打高尔夫球的人~1}{The person playing golf}]{%
  \label{sf:c}\includegraphics[width=0.4\textwidth]{golfer}}
\hspace{2em}
\subfigure[\bisubcaption{打高尔夫球的人~2}{The person playing golf}]{%
  \label{sf:d}\includegraphics[width=0.4\textwidth]{golfer}}
\end{minipage}
\bicaption[fig:bi]{}{双语子图题}{Fig.$\!$}{Bilingual subcaptions}
\end{figure}

引用：\ref{sf:c}、\ref{sf:d}。

\section{长子图题，考察换行与最后一行对齐}
\begin{figure}[!h]
\setlength{\subfigcapskip}{-1bp}
\centering
\subfigure[这是一个很长的子图题，用来把子图题撑到两行以上，看换行之后最后一行是居中还是左对齐，subcapcenterlast 选项管的就是这里]{\includegraphics[width=0.42\textwidth]{golfer}}
\hspace{1em}
\subfigure[短的]{\includegraphics[width=0.42\textwidth]{golfer}}
\bicaption[fig:long]{}{长子图题}{Fig.$\!$}{Long subcaptions}
\end{figure}

\section{子表}
\begin{table}[!h]
\centering
\subtable[左表]{\begin{tabular}{cc}\hline 甲 & 乙 \\\hline 1 & 2 \\\hline\end{tabular}}
\hspace{2em}
\subtable[右表]{\begin{tabular}{cc}\hline 丙 & 丁 \\\hline 3 & 4 \\\hline\end{tabular}}
\bicaption[tab:sub]{}{子表}{Tab.$\!$}{Subtables}
\end{table}

\section{三个子图排一行}
\begin{figure}[!h]
\centering
\subfigure[甲]{\includegraphics[width=0.28\textwidth]{golfer}}
\subfigure[乙]{\includegraphics[width=0.28\textwidth]{golfer}}
\subfigure[丙]{\includegraphics[width=0.28\textwidth]{golfer}}
\bicaption[fig:three]{}{三个子图}{Fig.$\!$}{Three subfigures}
\end{figure}

\section{最后一行居中（subcapcenterlast）}
% 这一节的效果要靠类选项 subcapcenterlast 打开，它在加载期决定，
% 正文里的 \hitsetup 改不动。用
%   OPT=subcapcenterlast bash tools/compile-subfigure.sh
% 编一遍与默认版对照。
\begin{figure}[!h]
\setlength{\subfigcapskip}{-1bp}
\centering
\subfigure[这是一个很长的子图题，用来把子图题撑到两行以上，看最后一行是不是居中，这一节把 subcapcenterlast 打开]{\includegraphics[width=0.42\textwidth]{golfer}}
\hspace{1em}
\subfigure[短的]{\includegraphics[width=0.42\textwidth]{golfer}}
\bicaption[fig:cl]{}{最后一行居中}{Fig.$\!$}{Center last line}
\end{figure}

\section{图目录里的子图条目}
见前置部分的插图索引。

\backmatter
\end{document}
